\paragraph{\trModuliSpace{}的\trOrientation。}\label{sub:orientations-moduli-spaces}
从现在起，我们回到惯常的框架：$(M,\varphi_M,\theta_M,I_M)$ 是一个%
\trCorneredBoundary{}的\trExactSymplectic\trManifold（仍带有某个 $\eta_M^2$），
并且我们只考虑\trExactLagrangianBrane——意思是底层\trLagrange\trSubmanifold{}
为\trExact{}的那些\trBrane。
设 $(L_0^{\#},L_1^{\#})$ 是一对\trExactLagrangianBrane，使得 $L_0$、$L_1$
\trTransverse{}相交。在每一交点 $y\in L_0\cap L_1$ 处的%
\trTangentSpace{}$(TM_y,\varphi_{M,y},I_{M,y},\eta^2_{M,y})$ 内，都有一对显然的%
\trLinear\trLagrange\trBrane
\begin{equation}\label{eq:linear-brane-pair}
\Lambda^{\#}_{k,y}=(\Lambda_{k,y}=(TL_k)_y,\ \alpha^{\#}_{k,y}=\alpha_k^{\#}(y),\
P^{\#}_{k,y}=(P_k^{\#})_y).
\end{equation}
取其在\eqref{eq:absolute-index-orientation-space} 中定义的\trIndex{}与%
\trOrientation{}空间，并记为 $i(y)=i(\Lambda^{\#}_{0,y},\Lambda^{\#}_{1,y})$、
$o(y)=o(\Lambda^{\#}_{0,y},\Lambda^{\#}_{1,y})$。更一般地，可以容许任意的
$(L_0,L_1)$，选取一组\trFloerDatum{}$(H,J)$，并考虑 $y\in\scr{C}(L_0,L_1)$。
在此情形，(12.1) 中的诸 $\Lambda^{\#}_{k,y}$ 位于\trApriori{}不同的\trVectorSpace{}
$TM_{y(k)}$ 中；不过，与\cref{sub:orientation-operators}中的讨论相平行，可以利用
$H$ 的\trHamiltonian{}“流”在 $y$ 附近的\trLinearization，连同
$Gr(TM_{y(t)})$ 上的诸\trSquaredPhase{}函数，把 $\Lambda^{\#}_{0,y}$ 移到
$TM_{y(1)}$。做完这一步之后，$i(y)$ 与 $o(y)$ 便可像从前一样定义。由
\eqref{eq:orientation-operator-index}，它们同样可以被看作是为
$(\Lambda^{\#}_{0,y},\Lambda^{\#}_{1,y})$ 任意选定的一个\trOrientationOperator{}
的\trIndex{}与\trDeterminantLine。我们将把这样的算子记作 $D_{H,y}$。

设 $S$ 是一个配备\trStriplikeEnd{}与\trBrane\trLabel{}（即\trLagrange\trLabel{}
连同\trBraneStructure）$L^{\#}_0,\ldots,L^{\#}_d$ 的 $(d{+}1)$ \trPointedDisc。
为诸对 $(L_0,L_d)$ 与 $(L_{k-1},L_k)$（$k=1,\ldots,d$）选取\trFloerDatum，并在
$S$ 上选取一个\trCompatible{}的\trPerturbationDatum。在\trRegular{}点
$u\in\scr{M}_S(y_0,\ldots,y_d)$ 处，有
\begin{equation}\label{eq:disc-dim-orientation}
\begin{gathered}
\dim_u(\scr{M}_S(y_0,\ldots,y_d))=i(y_0)-i(y_1)-\cdots-i(y_d),\\
(\lambda^{\mathrm{top}}T\scr{M}_S(y_0,\ldots,y_d))_u
\cong\lambda^{\mathrm{top}}\ker(D_{S,u})
\cong o(y_0)\otimes o(y_1)^{\vee}\otimes\cdots\otimes o(y_d)^{\vee}.
\end{gathered}
\end{equation}
用文字来说，第二行讲的是：该\trModuliSpace{}\emph{相对于诸端点处的\trOrientation{}
空间是\trCanonical{}地\trOrientation}的。\eqref{eq:disc-dim-orientation} 是%
\cref{thm:general-surface-index-formula} 的一个推论；不过与其援引那条命题，也许
更有启发性的做法是把相关的\trGluing{}论证重做一遍。选取\trOrientationOperator{}
$D_{H,y_k}$（$k=1,\ldots,d$）。按定义，这些是 $H$ 上的 $\bar\partial$-算子；
粗略地说，其\trBoundaryCondition{}的诸\trSubbundle{}上带有\trBraneStructure。
\trLinearization{}算子 $D_{S,u}$ 带有类似的数据。于是，把 $D_{S,u}$ 与
$D_{H,y_d},\ldots,D_{H,y_1}$ \trGluing{}起来（按此顺序），并与给定的诸%
\trBraneStructure{}\trCompatible，便得到 $(\Lambda^{\#}_{0,y_0},\Lambda^{\#}_{1,y_0})$
的一个\trOrientationOperator，我们称之为 $D_{H,y_0}$。\trGluing{}公式连同
\eqref{eq:orientation-operator-index} 立刻给出
\begin{equation}\label{eq:gluing-index-orientation}
\begin{gathered}
i(y_0)=\mathrm{index}(D_{S,u})+i(y_1)+\cdots+i(y_d),\\
o(y_0)\cong\det(D_{S,u})\otimes o(y_d)\otimes\cdots\otimes o(y_1).
\end{gathered}
\end{equation}
在\trRegular{}的情形，$D_{S,u}$ 是满的，且 $\ker(D_{S,u})\cong
T\scr{M}_S(y_0,\ldots,y_d)_u$，于是\eqref{eq:disc-dim-orientation} 随之得证。
\begin{remark}\label{thm:dual-pairing-remark}
说得更确切些，这里的约定是：对偶带有\trCanonical\trPairing{}
$o(y_k)^{\vee}\otimes o(y_k)\to\RR$。特别地，若从
\eqref{eq:disc-dim-orientation} 中的\trIsomorphism{}出发，与所有的 $o(y_k)$ 作%
\trTensorProduct{}（在右边、按递减次序），便重新得到
\eqref{eq:gluing-index-orientation}。
\end{remark}
我们需要知道，这些\trOrientation{}在把\trInhomogeneous\trPseudoholomorphicMap{}
\trGluing{}起来这一基本运算之下表现如何。取一个带\trBrane\trLabel{}
$(L^{\#}_0,\ldots,L^{\#}_n,L^{\#}_{n+m},\ldots,L^{\#}_d)$ 的 $(d{+}2{-}m)$
\trPointedDisc{}$S_1$，以及一个带\trLabel{}$(L^{\#}_n,\ldots,L^{\#}_{n+m})$ 的
$(m{+}1)$ \trPointedDisc{}$S_2$。把 $S_1$ 的第 $(n{+}1)$ 个出向端与 $S_2$ 的
入向端\trGluing{}起来（取某个\trGluing{}长度 $l$），便得到一个 $(d{+}1)$
\trPointedDisc{}$S$，它承袭\trLabel{}$(L^{\#}_0,\ldots,L^{\#}_d)$。为所有相关的%
\trLagrange\trSubmanifold{}对选取\trFloerDatum，并在 $S_1$、$S_2$ 上选取%
\trCompatible{}的\trPerturbationDatum，使之在 $S$ 上诱导出一个%
\trPerturbationDatum。现在设对某个 $\tilde{y}$，
$u_1\in\scr{M}_{S_1}(y_0,\ldots,y_n,\tilde{y},y_{n+m+1},\ldots,y_d)$ 与
$u_2\in\scr{M}_{S_2}(\tilde{y},y_{n+1},\ldots,y_{n+m})$ 是各自\trModuliSpace{}
中的\trRegular{}点。只要 $l$ 选得足够大，\trGluing{}过程便会给出一个\trRegular{}
的 $u\in\scr{M}_S(y_0,\ldots,y_d)$：先把诸 $u_k$ 朴素地拼合起来（借助%
\trPartitionOfUnity），再应用\trImplicitFunctionTheorem。这可以扩张到诸%
\trModuliSpace{}的邻近点，而且实际上给出了 $(u_1,u_2)$ 的邻域与 $u$ 的邻域之间
的一个\trLocalDiffeomorphism{}$g$。即便从这个含糊的描述出发也应当可以想见：
\trDifferential{}$Dg$ 逼近的是\trLinearization{}算子层面上的映射
\eqref{eq:delbar-gluing}。为了引出相关的结论，考察如下图式
\begin{equation}\label{eq:orientation-diagram}
\begin{tikzpicture}[line cap=round,line join=round]
\node (A) at (0.1,3.4) {$(\lambda^{\mathrm{top}}T\scr{M}_S)_u$};
\node (B) at (5.2,3.4)
 {$o(y_0)\otimes o(y_1)^{\vee}\otimes\cdots\otimes o(y_d)^{\vee}$};
\node (C) at (5.2,0.9) {$\displaystyle\begin{gathered}
o(y_0)\otimes\bigotimes_{k\le n}o(y_k)^{\vee}\otimes o(\tilde{y})^{\vee}
\otimes\bigotimes_{k>n+m}o(y_k)^{\vee}\\
\otimes\,o(\tilde{y})\otimes\bigotimes_{n<k\le n+m}o(y_k)^{\vee},\end{gathered}$};
\node (D) at (-0.75,0.35) {$\displaystyle\begin{gathered}
(\lambda^{\mathrm{top}}T\scr{M}_{S_1})_{u_1}\\
\otimes\,(\lambda^{\mathrm{top}}T\scr{M}_{S_2})_{u_2}\end{gathered}$};
\draw[->] (A) -- (B) node[midway,above]{$\cong$};
\draw[->] (C.north) -- (B.south);
\draw[->] (D) -- (C) node[pos=0.35,above]{$\cong$};
\draw[->] (-2.75,0.7) -- (-2.75,3.25)
 node[pos=0.72,right]{$(\lambda^{\mathrm{top}}Dg)_{(u_1,u_2)}$};
\end{tikzpicture}
\end{equation}
其中诸\trHorizontal\trIsomorphism{}即\eqref{eq:disc-dim-orientation}，而右侧
那个向上箭头是 $o(\tilde{y})$ 与其对偶之间的\trCanonical\trPairing。于是，
\eqref{eq:orientation-diagram} 在\trKoszul{}符号规则所规定的正负号之下交换，
该符号即 $(-1)^{\Delta}$，其中
$\Delta=\bigl(\sum_{k>n+m}i(y_k)\bigr)\cdot\mathrm{index}(D_{S_2,u_2})$。
这一符号的出现可以追溯到我们对\trDeterminantLine{}的定义，见
\cref{sub:determinant-lines}，特别是\eqref{eq:det-koszul-iso}。
上面的讨论可以向两个不同的方向推广。出于叙述上的考虑，我们对这两者分别论述，
把它们结合起来这一直截了当的任务则留给读者。
